\section{System}

\subsection{Interaction model}

\systemname{} represents an interaction at several linked levels. A
\emph{scenario step} captures a user-visible action, such as asking an exhibit
agent to describe a dinosaur skeleton. A \emph{capability use} identifies the
operation, the responsible provider, and the spatial targets. The capability is
connected through a \emph{runtime binding} to a concrete backend action. An
\emph{assertion} describes the outcome that should be observable, and a
\emph{validation case} links a stimulus to that expected outcome.

The required execution path is:

\begin{quote}
\texttt{ScenarioStep $\rightarrow$ CapabilityUse $\rightarrow$ Capability
$\rightarrow$ RuntimeBinding $\rightarrow$ RuntimeAction}.
\end{quote}

This path intentionally prevents a scenario step from invoking an endpoint
without an explicit agent capability. Agents additionally reference their
responsible zones and grounded assets. These relationships allow validators to
reject a capability use whose provider is not responsible for its spatial
target.

\subsection{Generation pipeline}

The pipeline starts from a structured room description. It normalizes scene
entities, derives scene semantics, proposes agent roles and capabilities,
materializes knowledge resources, and computes agent placement. A native
\systemname{} instance then connects these products to interaction scenarios and
runtime bindings. Deterministic validators check schemas and cross-reference
invariants before artifacts are made available to Unity or the backend.

The implementation separates probabilistic and deterministic stages. Language
models can propose semantic labels, roles, knowledge, and evaluation utterances.
Identifiers, geometric placement, relationship validation, artifact
materialization, and metric computation are deterministic. Each stage records
its inputs, configuration, outputs, and hashes so that a result can be
invalidated when a relevant upstream artifact changes.

\subsection{Unity/WebXR client}

The Unity client imports room assets and agent configurations, places embodied
agents in their assigned zones, and supports first-person and proximity-based
interaction. A visitor can communicate by text or speech. The client sends the
active project, agent, user utterance, and interaction context to the backend.
WebXR and the XR Interaction Toolkit provide the browser and immersive runtime,
while the same scene can be exercised in a desktop mode for deterministic
testing.

\subsection{Agent backend and handoffs}

The backend exposes setup, chat, speech, and project-management endpoints. Each
agent has a persona, knowledge scope, capabilities, and permitted handoff
targets. A request can be answered by the active agent, delegated to a more
appropriate agent, or rejected when no grounded target or responsible
capability exists. Runtime actions carry the interaction-model identifiers
needed to distinguish a successful HTTP request from a semantically correct
interaction.

\subsection{Runtime validation}

Events are written against the originating scenario, capability use, binding,
agent, and spatial target. Validation separates four questions:

\begin{enumerate}
  \item Was a syntactically valid runtime operation invoked?
  \item Was the intended spatial entity resolved?
  \item Was the responsible agent and capability selected?
  \item Did the observable response satisfy the case-specific assertion?
\end{enumerate}

This separation avoids treating network success or a single token overlap as
evidence of grounded interaction. It also provides a failure location that can
be surfaced in authoring tools or logs.

\subsection{Use of machine intelligence}

The system uses language models for semantic scene interpretation, agent-role
proposal, knowledge synthesis, conversational responses, and, in an explicitly
reported condition, semantic response assessment. Prompts, model identifiers,
sampling parameters, and repair attempts are versioned with the evaluation.
No language-model output is accepted as structurally valid without schema and
invariant checks. The Stanford Agentic Reviewer is not part of the system or the
scientific evaluation.
